% !PRINTED 339
\weberpage{339}
% running head: §. 91. | Invariants of the Curve of the Third Order. | 339

This equation takes a simpler form through the
substitution
\weq{11}{z = \frac{3m^2}{r^2}.}
It then becomes
\weq{12}{z^4+6Sz^2+Tz-3S^2=0.}
This biquadratic equation has the property that
its first invariant $A=0$, and that its discriminant
(Bd. I, \S.\,47):
\weq{13}{D_2 = -27(T^2+64S^3)^2 = -27D^2,}
is thus always negative. The quantity $D=T^2+64S^3$ takes for
the canonical form (5) the expression:
\weq{14}{h^3(h^3+8m^3)^3,}
and hence cannot vanish so long as the curve $f$ has no
singular point. (It is in fact nothing other than
the discriminant of the form $f$.)

From (13) it now follows that the biquadratic equation
 (12) has, for real curves $f$, a negative discriminant, and
consequently (Bd. I, \S.\,78) two real and two conjugate imaginary
roots, and since the product of all four roots, $-3S^2$,
is negative, \gesperrt{of the real roots one is positive},
one negative.

From $z$ one can calculate, with the help of (11) and the first equation
 (7), the two quantities $m$ and $r$. One obtains in this way
\weq{15}{m^3 = \frac{z^2}{9S+z^2}, \quad mr^2 = \frac{3z}{9S+z^2}, \quad r^6 = \frac{27z}{(9S+z^2)^2},}
and here, as one concludes from (12) by setting
$z^2=-9S$, the quantity $9S+z^2$ could vanish only if either
 $S$ or $D$ is zero. With this the coefficients $P$, $Q$, $R^3$
of the cubic equation \S.\,90, (14) are uniquely determined by $z$ and by
known quantities, and for the real positive root $z$
the values $m$ and $r$ become real, and consequently $y_1$, $y_2$, $y_3$ real as well.

A slight exception is formed by the case $S=0$, in which one
root of equation (12) vanishes, and the second real
root becomes $-\sqrt[3]{T}$, and for the corresponding values of $m$
and $r^6$ one obtains from (7) and (11)
\[
m = 0, 1; \quad r^6 = \frac{1}{T}, \ \frac{-27}{T};
\]
%CONT>>
